\textbf{Phần 1. Trắc nghiệm lựa chọn}
\setcounterpageref{ex}{0}

\begin{ex} %[2D4V20]
Tính giá trị phép tính: $38 + 27$.
\choice
{$64$}
{\True $65$}
{$66$}
{$55$}
\loigiai{
$38 + 27 = 65$.
\textbf{Đáp án của câu hỏi là B}
}
\end{ex}

\begin{ex} %[2D4V20]
Lớp 2A có $29$ học sinh, lớp 2B có $34$ học sinh. Cả hai lớp có bao nhiêu học sinh?
\choice
{$61$ học sinh}
{$62$ học sinh}
{$64$ học sinh}
{\True $63$ học sinh}
\loigiai{
$29 + 34 = 63$ (học sinh).
\textbf{Đáp án của câu hỏi là D}
}
\end{ex}

\begin{ex} %[2D4V20]
Tìm $x$ biết: $x - 36 = 28$.
\choice
{\True $64$}
{$62$}
{$63$}
{$65$}
\loigiai{
$x = 28 + 36 = 64$.
\textbf{Đáp án của câu hỏi là A}
}
\end{ex}

\begin{ex} %[2D4V20]
Tính nhanh: $25 + 18 + 35 + 12$.
\choice
{$80$}
{$85$}
{\True $90$}
{$95$}
\loigiai{
$(25 + 35) + (18 + 12) = 60 + 30 = 90$.
\textbf{Đáp án của câu hỏi là C}
}
\end{ex}

\begin{ex} %[2D4V20]
Một cửa hàng buổi sáng bán được $48\text{ kg}$ gạo, buổi chiều bán được $37\text{ kg}$ gạo. Cả ngày cửa hàng bán được bao nhiêu ki-lô-gam gạo?
\choice
{$83\text{ kg}$}
{\True $85\text{ kg}$}
{$84\text{ kg}$}
{$86\text{ kg}$}
\loigiai{
$48 + 37 = 85\text{ (kg)}$.
\textbf{Đáp án của câu hỏi là B}
}
\end{ex}

\begin{ex} %[2D4V20]
Cho phép tính: $\overline{a6} + 28 = 64$. Giá trị của chữ số $a$ là
\choice
{$2$}
{$4$}
{$1$}
{\True $3$}
\loigiai{
Ta có: $\overline{a6} = 64 - 28 = 36$. Do đó $a = 3$.
\textbf{Đáp án của câu hỏi là D}
}
\end{ex}

\begin{ex} %[2D4V20]
Một con đường dài $72\text{ m}$. Đội công nhân thứ nhất sửa được $35\text{ m}$, đội thứ hai sửa được $37\text{ m}$. Hai đội đã sửa được tất cả bao nhiêu mét đường?
\choice
{\True $72\text{ m}$}
{$70\text{ m}$}
{$71\text{ m}$}
{$73\text{ m}$}
\loigiai{
$35 + 37 = 72\text{ (m)}$. Hai đội đã sửa xong toàn bộ con đường.
\textbf{Đáp án của câu hỏi là A}
}
\end{ex}

\begin{ex} %[2D4V20]
Tìm tổng của số lớn nhất có hai chữ số khác nhau nhỏ hơn $50$ và số nhỏ nhất có hai chữ số giống nhau lớn hơn $20$.
\choice
{$70$}
{$72$}
{\True $71$}
{$69$}
\loigiai{
Số lớn nhất có 2 chữ số khác nhau nhỏ hơn 50 là 49. Số nhỏ nhất có 2 chữ số giống nhau lớn hơn 20 là 22.\nTổng: $49 + 22 = 71$.
\textbf{Đáp án của câu hỏi là C}
}
\end{ex}

\begin{ex} %[2D4V20]
Điền dấu thích hợp: $47 + 38 \dots 56 + 29$.
\choice
{$>$}
{\True $=$}
{$<$}
{Không so sánh được}
\loigiai{
$47 + 38 = 85$; $56 + 29 = 85$. Hai vế bằng nhau nên điền dấu $=$.
\textbf{Đáp án của câu hỏi là B}
}
\end{ex}

\begin{ex} %[2D4V20]
Một trang trại có $46$ con bò sữa và $38$ con bò thịt. Trang trại có tất cả bao nhiêu con bò?
\choice
{$82$ con bò}
{$83$ con bò}
{$85$ con bò}
{\True $84$ con bò}
\loigiai{
$46 + 38 = 84$ (con bò).
\textbf{Đáp án của câu hỏi là D}
}
\end{ex}

\begin{ex} %[2D4V20]
Đoạn dây thứ nhất dài $39\text{ cm}$, đoạn dây thứ hai dài $45\text{ cm}$. Nối hai đoạn dây lại (không kể nút nối) được đoạn dây dài bao nhiêu đề-xi-mét, xăng-ti-mét?
\choice
{\True $8\text{ dm } 4\text{ cm}$}
{$8\text{ dm } 3\text{ cm}$}
{$8\text{ dm } 5\text{ cm}$}
{$7\text{ dm } 4\text{ cm}$}
\loigiai{
Tổng chiều dài: $39 + 45 = 84\text{ cm} = 8\text{ dm } 4\text{ cm}$.
\textbf{Đáp án của câu hỏi là A}
}
\end{ex}

\begin{ex} %[2D4V20]
Tìm một số biết rằng lấy số đó trừ đi $29$ thì được $48$.
\choice
{$76$}
{$78$}
{\True $77$}
{$75$}
\loigiai{
Số đó là: $48 + 29 = 77$.
\textbf{Đáp án của câu hỏi là C}
}
\end{ex}

\textbf{Phần 2. Trắc nghiệm đúng sai}
\setcounterpageref{ex}{0}

\begin{ex} %[2D4V20]
Xét tính đúng sai của phép cộng có nhớ trong phạm vi 100:
\choiceTF
{\True Phép cộng $47 + 36$ có kết quả là $83$.}
{Tổng của hai số có hai chữ số luôn nhỏ hơn 90.}
{\True Khi cộng $28 + 37$, hàng đơn vị có $8 + 7 = 15$ viết 5 nhớ 1 sang hàng chục.}
{Nếu tổng bằng 75 và một số hạng là 38 thì số hạng kia là 47.}
\loigiai{
\begin{itemize}
    \item Ý a) Đúng: $47 + 36 = 83$.
    \item Ý b) Sai: Ví dụ: $50 + 50 = 100 > 90$.
    \item Ý c) Đúng: Quy tắc đặt tính rồi tính chuẩn.
    \item Ý d) Sai: Số hạng kia là: $75 - 38 = 37$ (không phải 47).
\end{itemize}
\begin{itemize}
    \item a) \textbf{Đúng}.
    \item b) \textbf{Sai}.
    \item c) \textbf{Đúng}.
    \item d) \textbf{Sai}.
\end{itemize}
}
\end{ex}

\begin{ex} %[2D4V20]
Hai kho gạo có tất cả 92 bao gạo, kho A có 48 bao gạo:
\choiceTF
{\True Kho B có 44 bao gạo.}
{\True Kho A nhiều hơn Kho B là 4 bao gạo.}
{Chuyển 4 bao gạo từ kho A sang kho B thì hai kho bằng nhau.}
{\True Nếu kho A nhập thêm 5 bao thì kho A có 53 bao gạo.}
\loigiai{
\begin{itemize}
    \item Ý a) Đúng: $92 - 48 = 44$ bao gạo.
    \item Ý b) Đúng: $48 - 44 = 4$ bao gạo.
    \item Ý c) Sai: Chỉ cần chuyển $4 : 2 = 2$ bao thì cả hai kho đều có 46 bao.
    \item Ý d) Đúng: $48 + 5 = 53$ bao. Đúng.
\end{itemize}
\begin{itemize}
    \item a) \textbf{Đúng}.
    \item b) \textbf{Đúng}.
    \item c) \textbf{Sai}.
    \item d) \textbf{Đúng}.
\end{itemize}
}
\end{ex}

\textbf{Phần 3. Trả lời ngắn}
\setcounterpageref{ex}{0}

\begin{ex} %[2D4V20]
Tính giá trị: $46 + 37$.
\shortans{83}
\loigiai{
\textbf{Đáp số:} $83$\\
$46 + 37 = 83$.
}
\end{ex}

\begin{ex} %[2D4V20]
Mẹ mua 28 kg gạo tẻ và 15 kg gạo nếp. Hỏi mẹ mua tất cả bao nhiêu ki-lô-gam gạo?
\shortans{43}
\loigiai{
\textbf{Đáp số:} $43$\\
$28 + 15 = 43\text{ kg}$.
}
\end{ex}

\begin{ex} %[2D4V20]
Tìm x biết: $x - 45 = 28$.
\shortans{73}
\loigiai{
\textbf{Đáp số:} $73$\\
$x = 28 + 45 = 73$.
}
\end{ex}

\begin{ex} %[2D4V20]
Tính nhanh: $27 + 34 + 13 + 16$.
\shortans{90}
\loigiai{
\textbf{Đáp số:} $90$\\
$(27 + 13) + (34 + 16) = 40 + 50 = 90$.
}
\end{ex}

\begin{ex} %[2D4V20]
Thùng thứ nhất chứa 38 lít dầu, thùng thứ hai chứa 47 lít dầu. Hỏi cả hai thùng chứa bao nhiêu lít dầu?
\shortans{85}
\loigiai{
\textbf{Đáp số:} $85$\\
$38 + 47 = 85\text{ L}$.
}
\end{ex}

\begin{ex} %[2D4V20]
Đoạn thẳng AB dài 36 cm, đoạn thẳng BC dài 48 cm. Tính độ dài đoạn thẳng AC biết ba điểm A, B, C thẳng hàng theo thứ tự đó.
\shortans{84}
\loigiai{
\textbf{Đáp số:} $84$\\
$36 + 48 = 84\text{ cm}$.
}
\end{ex}
